%! TEX root = ../am2-2018-2019.tex

\chapter{Ecuații și sisteme diferențiale}

\section{Ecuații rezolvabile prin cuadraturi}

Dacă ecuația are forma generală $y^{(n)} = f(x)$, atunci ea se rezolvă
prin $n$ integrări succesive și soluția generală depinde de $n$ constante
arbitrare.

Exemplu 1:
\[
  y^{\prime\prime} = \arcsin x + \frac{x}{\sqrt{1 - x^2}} - \frac{\pi}{6}.
\]

\emph{Soluție:} Integrăm succesiv și obținem:
\begin{align*}
  y' &= x\arcsin x - \frac{\pi}{6} x + c_1 \\
  y &= \frac{x^2}{2} \arcsin x + \frac{1}{4} x \sqrt{1 - x^2} - \frac{1}{4}\arcsin x - %
      \frac{\pi}{12}x^2 + c_1 x + c_2.
\end{align*}

\vspace{1cm}

Exemplu 2:
\[
  y^{\prime\prime\prime} = \sin x + \cos x.
\]
\emph{Indicație:} Ecuația se rezolvă prin integrări succesive.

%%%%%%%%%%%%%%%%%%%%%%%%%%%%%%%%%%%%%%%%%%%%%%%%%%%%%%%%%%%%%%%%%%%%%%

\section{Ecuații de forma $f(x, y^{(n)}) = 0$}

În această situație, distingem două cazuri:
\begin{enumerate}[(a)]
\item Dacă ecuația poate fi rezolvată în raport cu $y^{(n)}$, adică dacă
  $\frac{\partial f}{\partial y^{(n)}} \neq 0$, atunci obținem una sau mai multe
  ecuații ca în secțiunea anterioară;
\item Dacă ecuația $f(x, y^{(n)}) = 0$ nu este rezolvabilă cu ajutorul funcțiilor
  elementare în raport cu $y^{(n)}$, dar cunoaștem o reprezentare parametrică a
  curbei $f(u, v) = 0$, anume $u = g(t), v = h(t)$, cu $t \in [\alpha, \beta]$,
  atunci soluția generală poate fi obținută sub formă parametrică:
  \[
    x = g(t), \quad dy^{(n-1)} = y^{(n)} dx = h(t) g'(t) dt,
  \]
  de unde putem obține succesiv:
  \[
    \begin{cases}
      y^{(n-1)} &= h_1(t, c_1) \\
      \dotfill \\
      y(t) &= h_n(t, c_1, \dots, c_n).
    \end{cases}
  \]
\end{enumerate}

De exemplu:
\[
  x - e^{y^{\prime \prime}} + y^{\prime \prime} = 0.
\]

\emph{Soluție:} Putem nota $y^{\prime \prime} = t$ și obținem $x(t) = e^t - t$.
Deoarece $dy' = y^{\prime\prime} dx$, rezultă:
\[
  dy' = t(e^t - 1) dt \Rightarrow y' = -\frac{t^2}{2} + te^t - e^t + c_1.
\]
Dar, mai departe, $dy = y' dx$, deci:
\[
  dy = \big( - \frac{t^2}{2} + te^t - e^t + c_1)(e^t - 1) dt.
\]

În fine, soluția este:
\[
  y(t) = e^{2t}\big( \frac{t}{2} - \frac{3}{4} \big) + %
  e^t \big( -\frac{t^2}{2} + 1 + c_1 \big) + \frac{t^3}{6} - c_1 t + c_2.
\]

%%%%%%%%%%%%%%%%%%%%%%%%%%%%%%%%%%%%%%%%%%%%%%%%%%%%%%%%%%%%%%%%%%%%%%

\section{Ecuații de forma $f(y^{(n-1)}, y^{(n)}) = 0$}

Distingem două cazuri:

\begin{enumerate}[(a)]
\item Dacă ecuația este rezolvabilă prin funcții elementare în raport cu $y^{(n)}$,
  atunci notînd $z = y^{(n-1)}$, obținem $z' = f(z)$. Ecuația este cu variabile separabile
  și se rezolvă corespunzător, conducînd la $z = y^{(n-1)} = f_1(x, c_1)$, care este
  de tipul anterior;
\item Dacă ecuația nu este rezolvabilă prin funcții elementare în raport cu $y^{(n)}$,
  dar cunoaștem o reprezentare parametrică de forma:
  \[
    y^{(n-1)} = h(t), \quad y^{(n)} = g(t), \quad t \in [\alpha, \beta],
  \]
  atunci putem folosi relația $dy^{(n-1)} = y^{(n)}dx$, pentru a obține
  soluția sub formă parametrică:
  \begin{align*}
    x &= \int \frac{h'}{g} dt + c_1 \\
    y^{(n-1)} &= h(t) \\
    dy^{(n-2)} &= h(t) = \frac{hh'}{g} dt \\
    y^{(n-2)} &= \int \frac{hh'}{g} dt + c_2 \\
    &\vdots \\
    y &= \int y' dx + c_n.
  \end{align*}
\end{enumerate}

Exemplu: $y^{\prime \prime} = -\sqrt{1 - {y'}^2}$.

\emph{Soluție:} Facem schimbarea de funcție $z(x) = y'$ și ecuația diferențială
devine:
\[
  - \frac{dz}{\sqrt{1 - z^2}} = dx, \quad |z| < 1.
\]
Soluția generală este $\arccos z = x + c_1$, de unde $z = \cos(x + c_1)$.

Obținem mai departe $y(x) = \sin (x + c_1) + c_2$.

De asemenea, remarcăm și soluțiile particulare $y = \pm x + c$.

%%%%%%%%%%%%%%%%%%%%%%%%%%%%%%%%%%%%%%%%%%%%%%%%%%%%%%%%%%%%%%%%%%%%%%

\section{Ecuații diferențiale de forma $f(x, y^{(k)}, \dots, y^{(n)}) = 0$}

Ecuația se rezolvă făcînd schimbarea de funcție $y^{(k)} = z(x)$ și obținem
o ecuație de ordin $n - k$:
\[
  f(x, z, z', \dots, z^{(n-k)}) = 0,
\]
pe care o integrăm.

Exemplu: $(1 + x^2)y^{\prime \prime} = 2xy'$.

\emph{Soluție:} Cu schimbarea de funcție $y' = z(x)$, obținem ecuația:
\[
  \frac{z'}{z} = \frac{2x}{1 + x^2},
\]
iar apoi, prin integrare, avem $z = y' = c_1(1 + x^2)$. Rezultă, în fine:
\[
  y(x) = c_1 x + c_1 \frac{x^3}{3} + c_2.
\]

\vspace{1cm}

Exemplu: $y \cdot y^{\prime \prime} - 2{y^\prime}^2 = 0$.

\emph{Soluție:} Notăm $y' = z$ și obținem:
\[
  y(y' z + yz') - 2y^2 z^2 = y^2(z^2 + z') - 2y^2z^2 = 0,
\]
de unde avem $y = 0$ sau $z' - z^2 = 0$.

Din a doua variantă, deducem succesiv:
\[
  \frac{dz}{z^2} = dx \Rightarrow -\frac{1}{z} = x + c_1.
\]
Revenind la $y(x)$, avem:
\[
  -\frac{y}{y'} = x + c_1 \Rightarrow \frac{dy}{y} = -\frac{dx}{x + c_1} \Rightarrow %
  y(x) = \frac{c_2}{x + c_1}.
\]

%%%%%%%%%%%%%%%%%%%%%%%%%%%%%%%%%%%%%%%%%%%%%%%%%%%%%%%%%%%%%%%%%%%%%%

\section{Ecuații de forma $f(y', y^{\prime\prime}, \dots, y^{(n)}) = 0$}

Să remarcăm că aceste tipuri de ecuații nu conțin pe $y$. Așadar, putem lua $y'$
ca variabilă independentă și pe $y^{\prime\prime}$ ca funcție de $y'$.

De exemplu: $x^2 y^{\prime\prime} = {y^\prime}^2 - 2xy^\prime + 2x^2$.

\emph{Soluție:} Facem schimbarea de funcție $y' = z(x)$ și obținem o ecuație
Riccati:
\[
  z' = \frac{z^2}{x^2} - 2\frac{z}{x} + 2.
\]
Observăm soluția particulară $z_p = x$ și integrăm punînd $z = x + \frac{1}{u(x)}$.

Obținem succesiv:
\[
  z(x) = x + \frac{c_1 x}{x + c_1} \Rightarrow y'(x) = x + \frac{c_1 x}{x + c_1}.
\]
În fine, soluția este:
\[
  y(x) = \frac{x^2}{2} + c_1 x - c_1^2 \ln|x + c_1| + c_2.
\]

%%%%%%%%%%%%%%%%%%%%%%%%%%%%%%%%%%%%%%%%%%%%%%%%%%%%%%%%%%%%%%%%%%%%%%

\section{Ecuații autonome (ce nu conțin pe $x$)}

În cazul acestor ecuații, putem micșora ordinul cu o unitate dacă notăm $y' = p$
și luăm pe $y$ variabilă independentă.

\begin{remark}\label{rk:autonom-pierd}
  Există posibilitatea să pierdem soluții de forma $y = c$ prin această metodă,
  deci trebuie verificat ulterior dacă a fost cazul.
\end{remark}

Exemplu 1: $1 + {y'}^2 = 2yy'$.

\emph{Soluție:} Luăm $y' = p$ drept funcție și pe $y$ drept variabilă independentă.
Obținem:
\[
  y^{\prime\prime} = \frac{d}{dx} \big( \frac{dy}{dx} \big) = p \frac{dp}{dy}.
\]

Atunci ecuația devine $\frac{2pdp}{1 + p^2} = \frac{dy}{y}$, ce are drept soluție
$y = c_1(1 + p^2)$.

Acum trebuie să obținem pe $x$ ca funcție de $p$ și $c_1$. Deoarece $dx = \frac{1}{p}dy$,
iar $dy = 2c_1 pdp$, rezultă $dx = 2c_1 dp$. Așadar, $x = 2c_1 p + c_2$, iar soluția
generală este $x(p) = 2c_1 p + c_2$. Cum $y(p) = c_1(1 + p^2)$, rezultă:
\[
  y = c_1 + \frac{(x - c_2)^2}{4c_1}.
\]

\vspace{1cm}

Exemplu 2: $y^{\prime\prime} + {y'}^2 = 2e^{-y}$.

\emph{Soluție:} Notăm $y' = p$ și luăm ca necunoscută $p = p(y)$,
de unde obținem:
\[
  y^{\prime\prime} = \frac{dy'}{dx} = p'p \Rightarrow p'p + p^2 = 2e^{-y}.
\]
Dacă notăm $p^2 = z$, obținem o ecuație liniară neomogenă:
\[
  z' + 2z = 4e^{-y},
\]
ce are ca soluție generală $z(y) = c_1 e^{-2y} + 4e^{-y}$. Revenind la $y$,
avem:
\[
  z = p^2 = {y'}^2 \Rightarrow y' = \pm \sqrt{c_1 e^{-2y} + 4e^{-y}},
\]
adică ecuația:
\[
  \frac{dy}{\pm\sqrt{c_1 e^{2y} + 4e^{-y}}} = dx,
\]
ce are ca soluție: $x + c_2 = \pm \frac{1}{2} \sqrt{c_1 + 4e^y}$. Echivalent, în forma
implicită:
\[
  e^y + \frac{c_1}{4} = (x + c_2)^2.
\]

\vspace{1cm}

Similar putem proceda și în cazul:
\[
  y^{\prime\prime\prime} + y^{\prime\prime} = \sin x.
\]

\emph{Indicație:} Putem nota $y^{\prime\prime} = z$ și atunci ecuația devine:
\[
  z' + z = \sin x,
\]
care este o ecuație liniară neomogenă, cu soluția:
\[
  z(x) = e^{-x}(c_1 + \frac{e^x}{2}(\sin x - \cos x)).
\]
Mai departe, integrăm succesiv:
\begin{align*}
  y'(x) &= \int z(x) dx + c_2 = -c_1 e^{-x} - \frac{1}{2} \cos x - \frac{1}{2}\sin x + c_2 \\
  y(x) &= \int y'(x) dx + c_3 = c_1 e^{-x} - \frac{1}{2} \sin x + \frac{1}{2} \cos x + c_2 x + c_3.
\end{align*}

%%%%%%%%%%%%%%%%%%%%%%%%%%%%%%%%%%%%%%%%%%%%%%%%%%%%%%%%%%%%%%%%%%%%%%

\section{Ecuații liniare de ordinul $n$ cu coeficienți constanți}

Forma generală este:
\[
  L[y] = a_0 y^{(n)} + a_1 y^{(n-1)} + \dots + a_n y = f(x),
\]
cu $a_i \in \RR$ constante.

Dacă avem $f(x) = 0$, atunci ecuația se numește \emph{omogenă}.

Avem o metodă algebrică de a rezolva această ecuație, folosind:

\begin{definition}\label{def:pol-car-ec}
  Se numește \emph{polinomul caracteristic} atașat ecuației omogene $L[y] = 0$
  polinomul:
  \[
    F(r) = a_0 r^n + a_1 r^{n-1} + \dots + a_n,
  \]
  iar $F(r) = 0$ se numește \emph{ecuația caracteristică} atașată ecuației diferențiale.
\end{definition}

Folosind această noțiune, putem rezolva direct ecuația, ținînd seama de următoarele
cazuri posibile:

\begin{theorem}\label{thm:ec-car-dif}
  \begin{enumerate}[(1)]
  \item Dacă ecuația caracteristică $F(r) = 0$ are rădăcini reale și distincte $r_i$,
    atunci un sistem fundamental de soluții este dat de:
    \[
      \big\{ y_i(x) = e^{r_i x} \big\}_{i=1,\dots,n}.
    \]
  \item Dacă printre rădăcinile lui $F(r)$ există și rădăcini multiple, de exemplu
    $r_1$, cu ordinul de multiplicitate $p$, atunci pentru această rădăcină avem
    $p$ soluții liniar independente (pe lîngă celelalte):
    \[
      y_1(x) = e^{r_1 x}, y_2(x) = xe^{r_1 x}, \dots, y_p(x) = x^{p-1} e^{r_1 x}.
    \]
  \item Dacă printre rădăcinile ecuației caracteristice avem și rădăcini complexe,
    de exemplu $r = a + ib$ și $\overline{r} = a - ib$, atunci fiecărei perechi
    de rădăcini complexe conjugate îi corespund două soluții liniar independente
    (pe lîngă celelalte):
    \[
      y_1(x) = e^{ax} \cos bx, \quad y_2(x) = e^{ax} \sin bx.
    \]
  \item Dacă ecuația caracteristică are rădăcini complexe ca mai sus, cu ordinul
    de multiplicitate $p$, atunci lor le corespund $2p$ soluții liniar independente:
    \[
      \big\{ y_j(x) = x^{j-1} e^{ax} \cos bx \big\}_{j=1, \dots p}, \quad
      \big\{ y_k(x) = x^{k-p-1} e^{ax} \sin bx\big\}_{k=p+1, \dots, 2p}.
    \]
  \end{enumerate}
\end{theorem}

De exemplu: $y^{\prime\prime} - y = 0$, cu condițiile $y(0) = 2$ și $y'(0) = 0$.

\emph{Soluție:} Ecuația caracteristică este $r^2 - 1 = 0$, deci $r_{1,2} = \pm 1$.
Sîntem în primul caz al teoremei, deci un sistem fundamental de soluții este:
\[
  y_1(x) = e^{-x}, \quad, y_2(x) = e^x.
\]
Soluția generală este $y(x) = c_1 e^{-x} + c_2 e^x$.

Folosind condițiile Cauchy, obținem $c_1 = c_2 = 1$, deci soluția particulară este:
\[
  y(x) = e^{-x} + e^x.
\]

\vspace{1cm}

\begin{remark}\label{rk:neomogen}
  Dacă ecuația inițială $L[y] = f(x)$ \emph{nu} este omogenă, putem rezolva folosind
  \emph{metoda variației constantelor (Lagrange)}.
\end{remark}

În acest caz, metoda variației constantelor presupune următoarele etape.
Fie ecuația neomogenă scrisă în forma:
\[
  a_0(x) y^{(n)} + a_1(x) y^{(n - 1)} + \dots + a_{n-1}(x) y' + a_n(x) y = f(x).
\]

Pentru simplitate, vom presupune $n = 2$ și fie o soluție particulară de forma:
\[
  y_p(x) = c_1(x) y_1 + c_2(x) y_2.
\]

Atunci, prin metoda variației constantelor, vom determina funcțiile $c_1(x), c_2(x)$
ca soluții ale sistemului algebric:
\[
  \begin{cases}
    c'_1(x) y_1 + c'_2(x) y_2 &= 0 \\
    c'_1(x) y'_1 + c'_2(x) y'_2 &= \frac{f(x)}{a_0(x)}
  \end{cases}
\]

\vspace{1cm}

De exemplu: $y^{\prime \prime} + 3y' + 2y = \frac{1}{e^x + 1}$.

\emph{Soluție:} Asociem ecuația omogenă, care are ecuația caracteristică
$r^2 + 3r + 2 = 0$, cu rădăcinile $r_1 = -1, r_2 = -2$. Așadar, soluția generală
a ecuației omogene este:
\[
  \overline{y}(x) = c_1 e^{-x} + c_2 e^{-2x}.
\]

Determinăm o soluție particulară $y_p(x)$ cu ajutorul metodei variației constantelor.
Mai precis, căutăm:
\[
  y_p(x) = c_1(x) e^{-x} + c_2(x) e^{-2x}.
\]

Înlocuind în ecuația neomogenă, rezultă sistemul:
\[
  \begin{cases}
    c_1'(x) e^{-x} + c_2'(x) e^{-2x} &= 0 \\
    -c_1'(x) e^{-x} - 2c_2'(x) e^{-2x} &= \frac{1}{1 + e^x}
  \end{cases}
\]

Rezolvînd ca pe un sistem algebric, obținem:
\[
  c_1'(x) = \frac{e^x}{1 + e^x}, \quad c_2'(x) = -\frac{e^{2x}}{1 + e^x},
\]
de unde obținem:
\[
  c_1(x) = \ln(1 + e^x), \quad c_2(x) = -e^x + \ln(1 + e^x).
\]

În fine, înlocuim în soluția particulară $y_p$ și apoi în cea generală,
$y(x) = \overline{y}(x) + y_p(x)$.

\vspace{1cm}

Un alt exemplu: $y^{\prime\prime} - y = 4e^x$.

\emph{Soluție:} Ecuația caracteristică $r^2 - 1$ are rădăcinile $r_{1,2} = \pm 1$,
deci soluția generală a ecuației liniare omogene este $y_0(x) = c_1 e^x + c_2e^{-x}$.

Căutăm acum o soluție particulară a ecuației neomogene, folosind metoda variației
constantelor. Așadar, căutăm $y_p(x) = c_1(x) e^x + c_2(x) e^{-x}$. Din condiția
ca $y_p$ să verifice ecuația liniară neomogenă, obținem sistemul:
\[
  \begin{cases}
    c_1'(x) e^x + c_2'(x) e^{-x} = 0 \\
    c_1'(x) e^x - c_2'(x) e^{-x} = 4e^x
  \end{cases}
\]
Rezultă $c'_1(x)e^x = 2e^x$, deci $c'_1(x) = 2 \Rightarrow c_1(x) = 2x$, iar
$c'_2(x) e^{-x} = -2e^x$, deci $c_2(x) = -e^{2x}$.

În fine, soluția particulară este $y_p(x) = 2xe^x - e^x$, iar soluția generală:
\[
  y(x) = y_0(x) + y_p(x) = c_1e^x + c_2 e^{-x} + e^x(2x - 1).
\]


%%%%%%%%%%%%%%%%%%%%%%%%%%%%%%%%%%%%%%%%%%%%%%%%%%%%%%%%%%%%%%%%%%%%%%

\section{Metoda eliminării pentru sisteme diferențiale liniare}

Putem reduce sistemele diferențiale de ordinul I la ecuații de ordin superior.

De exemplu, să pornim cu sistemul:
\[
  \begin{cases}
    x' + 5x + y &= 7e^t - 27 \\
    y' - 2x + 3y &= -3e^t + 12
  \end{cases}
\]

\emph{Soluție:} Din prima ecuație, scoatem $y$ și derivăm:
\[
  y = 7e^t - 27 - 5x - x' \Rightarrow y' = 7e^t - 5 - x^{\prime\prime}.
\]
Înlocuim în a doua ecuație și obținem o ecuație liniară de ordin superior:
\[
  x^{\prime\prime} + 8x' + 17x = 31e^t - 93.
\]
Ecuația caracteristică este $r^2 + 8r + 17 = 0$, cu rădăcinile $r_{1,2} = -4 \pm i$.
Atunci soluția generală a ecuației omogene este:
\[
  \overline{x}(t) = e^{-4t}(c_1 \cos t + c_2 \sin t).
\]
Mai departe, căutăm o soluție particulară a ecuației neomogene folosind metoda
variației constantelor, a lui Lagrange:
\[
  x_p(t) = e^{-4t}(c_1(t) \cos t + c_2(t) \sin t).
\]
Determinăm funcțiile $c_1(t), c_2(t)$ din sistemul:
\[
  \begin{cases}
    c'_1(t) \cos t e^{-4t} + c_2'(t) \sin t e^{-4t} &= 0 \\
    c'_1(t) (-\sin t e^{-4t} - 4\cos t e^{-4t}) + %
    c'_2(t) (\cos t e^{-4t} - 4\sin t e^{-4t}) &= 31e^t - 93.
  \end{cases}
\]
Rezultă:
\begin{align*}
  c'_1(t) &= -31 \sin t e^{5t} + 93 \sin t e^{4t} \\
  \Rightarrow c_1(t) &= -\frac{31}{26} e^{5t}(5 \sin t - \cos t) + %
                        \frac{93}{17} e^{4t}(4 \sin t - \cos t) \\
  c'_2(t) &= 31 \cos t e^{5t} - 93\cos t e^{4t} \\
  \Rightarrow c_2(t) &= \frac{31}{26} e^{5t}(5 \sin t + \cos t) - %
                       \frac{93}{17} e^{4t}(4 \sin t + \cos t).
\end{align*}
În fine, obținem:
\[
  x(t) = e^{-4t}(c_1 \cos t + c_2 \sin t) + \frac{31}{26}e^t - \frac{93}{17},
\]
iar mai departe:
\[
  y(t) = e^{-4t}[(c_1 - c_2) \sin t - (c_1(t) - c_2) \cos t) - \frac{2}{13} e^t + \frac{6}{17}.
\]

\vspace{1cm}

Alt exemplu:
\[
  \begin{cases}
    y'_1 &= -y_2 + 1 \\
    x^2 y'_2 &= -2y_1 + x^2 \ln x
  \end{cases}
\]

\emph{Soluție:} Derivăm prima ecuație și obținem $y'_2 = -y_1^{\prime\prime}$.
Înlocuim în a doua ecuație și obținem:
\[
  x^2 y_1^{\prime\prime} - 2y_1 = -x^2 \ln x,
\]
care este o ecuație Euler. Îi asociem o ecuație omogenă $x^2y_1^{\prime\prime} - 2y_1 = 0$.
Facem schimbarea $x = e^t$ și obținem $y_1^{\prime\prime} - y'_1 - 2y_1 = 0$.
Ecuația caracteristică asociată este $r^2 - r - 2 = 0$, cu rădăcinile $r_1 = 2$
și $r_2 = -1$. Atunci soluția generală este:
\[
  \begin{cases}
    \overline{y}_1(t) &= c_1 e^{2t} + c_2 e^{-t} \\
    \overline{y}_1(x) &= c_1 x^2 + c_2 \frac{1}{x}.
  \end{cases}
\]

Căutăm o soluție particulară a ecuației neomogene cu metoda variației constantelor,
deci:
\[
  y_{1p}(x) = c_1(x) x^2 + c_2(x) \frac{1}{x}.
\]
Determinăm funcțiile $c_1(x), c_2(x)$ din sistemul:
\[
  \begin{cases}
    c'_1(x) x^2 + c'_2(x) \frac{1}{x} &= 0 \\
    2xc'_1(x) - c'_2(x) \frac{1}{x^2} &= \frac{-x^2 \ln x}{x^2}
  \end{cases}
\]

Atunci $c'_1(x) = -\frac{\ln x}{3x^3}$, iar $c'_2(x) = \frac{\ln x}{3}$.
Așadar:
\begin{align*}
  c_1(x) &= \frac{1}{6x^2} \ln x + \frac{1}{6x} \ln x + \frac{1}{6x} \\
  c_2(x) &= \frac{1}{3} (x \ln x - x).
\end{align*}

Rezultă:
\begin{align*}
  y_{1p}(x) &= \frac{1}{6} x \ln x + \frac{1}{2} \ln x + \frac{x}{6} - \frac{1}{3} \\
  y_1(x) &= c_1 x^2 + c_2 \frac{1}{x} + \frac{1}{6} x \ln x + \frac{1}{2} \ln x + %
           \frac{x}{6} - \frac{1}{3} \\
  y_2(x) &= 1 - y'_1 = \frac{1}{6} \ln x - 2c_1 x - c_2 \frac{1}{x^2} + \frac{1}{2x} - \frac{2}{3}.
\end{align*}

